\documentclass[11pt]{article}
\usepackage[utf8]{inputenc} % Pour les caractères accentués
\usepackage[french]{babel} % Pour la langue française
\usepackage{hyperref} % Pour les liens cliquables dans le PDF
\usepackage{titlesec} % Pour personnaliser les titres
\usepackage{graphicx}
\usepackage{eurosym}
\usepackage{amsmath}
\usepackage{amssymb}
\usepackage{enumitem}
\usepackage{array}
\usepackage{bbold}
\usepackage{hyperref}
\usepackage[T1]{fontenc}



% Supprime "Chapitre" et garde juste le numéro
\titleformat{\chapter}[hang]
    {\normalfont\huge\bfseries}
    {\thechapter.}
    {0.5em}
    {}
    

\setlength{\parindent}{0pt}

\title{Projet - Fondements mathématiques pour l'aide à la décision}
\author{Lucie Boucheron, Feryel Benameur, Clara Baignères}
\date{}
\begin{document}
\maketitle

\tableofcontents


\section{Introduction}

Dans le cadre de notre cours « Fondements mathématiques pour l’aide à la décision », nous devons réaliser un projet qui a pour objectif d’étudier la polarisation des préférences exprimées lors de processus électoraux. 

Nous avons un ensemble de votantes et de candidates, ainsi que différents processus de votes. Une élection est polarisée lorsque l’électorat peut être divisé en deux clusters de votantes distincts. Pour cela, différentes techniques de distance et de clusterisation sont utilisées dans notre projet, telles que les distances de Hamming et de Spearman, ainsi que la clusterisation K-Means. L’ensemble du code Python a été réalisé sous VS Code grâce à un notebook Jupyter. Le travail de groupe a pu être effectué via GitHub et l’utilisation de git grâce à un système de branches.

\section{Répartition du travail}

Ce travail qui avait pour vocation de comprendre la polarisation des préférences exprimées lors de processus électoraux à été réalisé en plusieurs étapes. 

Dans un premier temps, Lucie Boucheron a résolu les questions 1 à 6 (incluses). Elle a pu créer les méthodes principales pour générer des profils aléatoires (en contrôlant la polarisation de l'échantillon généré), nos types de votes (par approbation ou par ordre totaux) et générer les profils en fonction du type de vote donné. Elle a aussi pu nous éclairer sur les calculs des axiomes et leur compréhension pour le projet.

Parallèlement, Féryel Benameur s’est occupée des questions 7 à 10. L’objectif était de comprendre les principes des fonctions de distances de Hamming et de Spearman de façon mathématiques en s’aidant notamment des axiomes vérifiés par Lucie. Dans un second temps, il était question d’implémenter en Python les formules données par ces fonctions de distance.

Enfin, dans un dernier temps, Clara Baignères s’est occupée des questions 11 à 13 comprises. Elle a pu implémenter les méthodes pour calculer $u_1^*$ et $u_2^*$ en fonction du type de vote donné. Elle s’est aussi occupée de l’implémentation de l’algorithme des K-Means avec $k=2$ dans le cas de vote par approbation. Ces méthodes et calculs ont été d’une grande importance et ont pu être réutilisés en appelant simplement les fonctions pour les questions 14 et 15 réalisées par Féryel. Clara s'est également chargée de la mise en page en \LaTeX~de ce rapport, grâce aux contributions des trois participantes.

\section{Réponses aux questions}

\subsection{Question 1}

Notre génération de profils pour un vote par approbation repose sur la polarisation choisie en argument de notre fonction. 

La méthode commence par générer un vote $a$ aléatoire. Chaque votante se verra attribuer $a$ ou $\bar{a}$. Pour une polarisation de 0, toutes les votantes reçoivent $a$. Pour une polarisation de 1, chaque votante se voit attribuer $a$ ou $\bar{a}$ avec la même probabilité. On voit que dans le cas de polarisation 0, l'échantillon sera \emph{exactement} de polarisation 0. Dans le cas de polarisation 1, il sera de polarisation 1 en moyenne.

Pour une polarisation comprise entre 0 et 1, on biaise simplement les probabilités.
	
Enfin, pour apporter plus de réalisme, notre fonction accepte en paramètre un bruit (potentiellement nul) qui viendra légèrement modifier le vote attribué à chaque participante.

On a alors bien l'idée que plus la polarisation est faible, plus les bulletins seront similaires, et plus la polarisation est grande, plus les bulletins peuvent se répartir en deux groupes. 

\subsection{Question 2}

Pour les votes par ordres totaux, le principe est le même. La modification par le bruit consiste à permuter deux candidates dans le vote. 
 
\subsection{Question 4}
 
\subsubsection{Axiome 1 : régularité}

Le premier axiome n'est pas respecté. Prenons le premier point à vérifier, soit $\varphi(p) = 0$ ssi $p = p_a$. Nous allons montrer qu'il existe un $a$ tel que $\varphi(p_a) \neq 0$.

Si on prend $a = \{0\}^m$, on a alors $n_{c_k,c_l}(p) = 0$ car il n'y a aucun vote $a$ où $a_v[k] = 1 \land a_v[l]=0$, puisque pour tous les votes, l'ensemble des bulletins est à $0$.

On a alors que $d_{c_k,c_l}(p) = 0$ et donc :

\[
\varphi(p) = \sum_{\{c_k,c_l\} \in C^2} 
\frac{n - d_{c_k,c_l}(p)}{n \binom{m}{2}} 
= \sum_{\{c_k,c_l\} \in C^2} \frac{1}{\binom{m}{2}} 
\]

Or il y a exactement $\binom{m}{2}$ couples distincts $({c_k,c_l})$. Donc : 

\[
\varphi^2(p_a) = \binom m2 \times \frac {1}{\binom m2} = 1
\]

De manière générale, cette partie de l'axiome (au moins) n'est jamais vérifiée quand tous les bulletins de tous les votes sont égaux (que des 1 ou que des 0).


\subsubsection{Axiome 2 : neutralité}

Soit $\pi$ une permutation de $\{1, \ldots, m\}$.

\begin{align*}
\varphi^2(p^\pi) 
&= \sum_{\{c_k,c_l\}} 
\frac{n - d_{c_k,c_l}(p^\pi)}{n \binom{m}{2}} \\
&= \sum_{\{c_k,c_l\}}
\frac{n - d_{c_{\pi^{-1}(k)}, c_{\pi^{-1}(l)}}(p)}{n \binom{m}{2}} \\
&= \sum_{\{c_k,c_l\}}
\frac{n - d_{c_k,c_l}(p)}{n \binom{m}{2}} \\
&= \varphi^2(p)
\end{align*}

Donc le second axiome est vérifié.


\subsubsection{Axiome 3 : anonymité}

Le principe est le même que pour l'axiome 2. Le fait de permuter les lignes (donc les votantes) ne change pas les valeurs de $n_{c_k,c_l}$. On a donc que $n_{c_k,c_l}(p) = n_{c_k,c_l}({}^\sigma p)$, et ainsi $\varphi^2(p) = \varphi^2({}^\sigma p)$.\\

Donc le troisième axiome est vérifié.


\subsubsection{Axiome 4 : invariance par réplication}

On note $kp$ l'ensemble de $k$ copies de notre profil. Alors :
\[
n_{c_k,c_l}(kp) = k n_{c_k,c_l}(p)
\quad \text{et} \quad
n_{c_l,c_k}(kp) = k n_{c_l,c_k}(p)
\]

Donc :
\[
d_{c_k,c_l}(kp) = k d_{c_k,c_l}(p)
\quad \text{et} \quad
kn - d_{c_k,c_l}(kp) = k(n - d_{c_k,c_l}(p))
\]

Ainsi :
\begin{align*}
\varphi^2(kp) 
&= \sum_{\{c_k,c_l\} \in C^2}
\frac{k(n - d_{c_k,c_l}(p))}{kn \binom{m}{2}} \\
&= \sum_{\{c_k,c_l\} \in C^2}
\frac{n - d_{c_k,c_l}(p)}{n \binom{m}{2}} \\
&= \varphi^2(p)
\end{align*}

Donc l’axiome 4 est vérifié.

\subsection{Question 6}

Pour ne pas déterminer la représentation graphique seulement à partir d'un profil par polarisation, on boucle sur toutes les valeurs de polarisation possibles et on génère plusieurs profils pour lesquels on calcule $\varphi^2$. On calcule enfin la moyenne des résultats. 

On remarque que, pour le vote par approbation comme pour le vote par ordre total, la fonction $\varphi^2$ se rapproche fortement d'une fonction linéaire $y=x$ ce qui est logique puisque $\varphi^2$ est sensé représenter notre polarisation. 

\subsection{Question 7 (pour la distance de Hamming)}

\subsubsection{Propriété 1 : positivité}

Par définition 
\[
d_H(a_{v_k}, a_{v_l}) = \sum_{i=1}^n \mathbb{1}_{a_{v_k}[i] \neq a_{v_l}[i]},
\]
La somme est positive comme somme d'éléments positifs ou nuls.  
D’où $d_H(a_{v_k},a_{v_l}) \geq 0$.
\\

Donc la première propriété est vérifiée.


\subsubsection{Propriété 2~: séparation}

\begin{align*}
d(a_{v_k}, a_{v_l}) = 0
&\Leftrightarrow
\sum_{i=1}^m \mathbb{1}_{a_{v_k}[i] \neq a_{v_l}[i]} = 0 \\
&\Leftrightarrow
\mathbb{1}_{a_{v_k}[i] \neq a_{v_l}[i]} = 0 \quad \forall i \in [\![1, \ldots, m]\!] \\
&\Leftrightarrow
a_{v_k}[i] = a_{v_l}[i] \quad \forall i \in [\![1, \ldots, m]\!]
\end{align*}
\\

Donc la seconde propriété est vérifiée.


\subsubsection{Propriété 3~: symétrie}

Pour tout $a_{v_k}, a_{v_l} \in \mathcal{A}$, on a :
\[
d_H(a_{v_k}, a_{v_l}) = d_H(a_{v_l}, a_{v_k})
\Leftrightarrow
\sum_{i=1}^m \mathbb{1}_{a_{v_k}[i] \neq a_{v_l}[i]} = 
\sum_{i=1}^m \mathbb{1}_{a_{v_l}[i] \neq a_{v_k}[i]}.
\]

Cela est vrai car l’inégalité ainsi que la somme sont commutatives donc changer l'ordre des termes n'impacte pas l'opération.
\\

Donc la troisième propriété est vérifiée. 

\subsubsection{Propriété 4 : inégalité triangulaire}

Soient trois bulletins $a = a_{v_k}$, $b = a_{v_l}$ et $c = a_{v_z}$ dans $\mathcal{A}$.

On a donc :
\[
d_H(a,c) = \sum_{i=1}^m \mathbb{1}_{a[i] \neq c[i]},
\quad
d_H(a,b) = \sum_{i=1}^m \mathbb{1}_{a[i] \neq b[i]},
\quad
d_H(b,c) = \sum_{i=1}^m \mathbb{1}_{b[i] \neq c[i]}.
\]

Or, pour tout indice $i \in \{1,\ldots,m\}$, on a~:

\[
\mathbb{1}_{a[i] \neq c[i]}
\le 
\mathbb{1}_{a[i] \neq b[i]} + \mathbb{1}_{b[i] \neq c[i]}.
\]

En effet :

\begin{itemize}
  \item Si $a[i] = c[i]$, alors $\mathbb{1}_{a[i] \neq c[i]} = 0$. Comme le membre de droite est positif ou nul, l'inégalité est vérifiée.
  \item Si $a[i] \neq c[i]$, alors $a[i]$, $b[i]$, $c[i]$ ne peuvent pas tous être différents simultanément (car ils valent 0 ou 1).  
  Ainsi, $b[i]$ diffère au moins de l’un des deux, et donc au moins l’une des deux indicatrices à droite vaut 1 ; comme le terme de gauche est inférieur ou égal à 1, l'inégalité est vérifiée.
\end{itemize}


En sommant sur toutes les coordonnées $i \in \{1,\ldots,m\}$, on obtient :
\[
\sum_{i=1}^m \mathbb{1}_{a[i] \neq c[i]} 
\le
\sum_{i=1}^m 
\big(
\mathbb{1}_{a[i] \neq b[i]} + \mathbb{1}_{b[i] \neq c[i]}
\big).
\]

Par la linéarité de la somme on a donc :
\[
d_H(a,c) \le d_H(a,b) + d_H(b,c).
\]
\\

La quatrième propriété est donc vérifiée.


\subsection{Question 7 (pour la distance de Spearman)}

\subsubsection{Propriété 1 : positivité}

Soient $a = \preceq_{v_k}$, $b = \preceq_{v_l}$, avec $a,b \in \mathcal{L}$.

\[
d_S(a,b) = \sum_{i=1}^m |r_a(i) - r_b(i)|
\]
avec $r_a(i)$ le rang de la candidate $i$ dans le vote de la votante $a$.  
Ce résultat est donc toujours positif par définition de la valeur absolue.
\\

Donc la première propriété est vérifiée.


\subsubsection{Propriété 2 : séparation}

On a~:
\begin{align*}
d_S(a,b) = 0 
&\Leftrightarrow 
\sum_{i = 1}^m |r_a(i) - r_b(i)| = 0\\
&\Leftrightarrow 
|r_a(i) - r_b(i)| = 0 \quad \forall i \in [\![1, \ldots, m]\!]\\
&\Leftrightarrow 
r_a(i) = r_b(i) \quad \forall i \in [\![1, \ldots, m]\!]
\end{align*}


Donc la seconde propriété est vérifiée.


\subsubsection{Propriété 3 : symétrie}

Pour tout $a,b \in \mathcal{L}$~:
\begin{align*}
d_S(a,b) 
&= 
\sum_{i = 1}^m |r_a(i) - r_b(i)|\\
&= 
\sum_{i = 1}^m |r_b(i) - r_a(i)|\\
&= 
d_S(b,a) 
\end{align*}

Donc la troisième propriété est vérifiée.


\subsubsection{Propriété 4 : inégalité triangulaire}

\begin{align*}
d_S(a,b) 
&= \sum_{i=1}^m |r_a(i) - r_b(i)| \\
&= \sum_{i=1}^m |r_a(i) - r_c(i) + r_c(i) - r_b(i)| \\
&\leq \sum_{i=1}^m \left(|r_a(i) - r_c(i)| + |r_c(i) - r_b(i)|\right) \text{ par l'inégalité triangulaire de la valeur absolue}\\
&= \sum_{i=1}^m |r_a(i) - r_c(i)| + \sum_{i=1}^m |r_c(i) - r_b(i)| \text{ par linéarité de la somme} \\
&= d_S(a,c) + d_S(c,b)
\end{align*}

On a donc bien que :
\[
d_S(a,b) \leq d_S(a,c) + d_S(c,b).
\]
\\


Donc la quatrième proposition est vérifiée.


\subsection{Question 9}

\subsubsection{Axiome 1 : régularité}

%On souhaite montrer que 
%\[
%\varphi_{d_H}(p) \in [0,1]
%\quad \text{et que} \quad
%\varphi_{d_H}(p) = 0 \text{ ssi } p = p_a
%\quad \text{et} \quad
%\varphi_{d_H}(p) = 1 \text{ ssi } p = p_{a, \bar{a}}.
%\]

\textbf{On montre que $\varphi_{d_H}(p) = 0$ ssi $p = p_a$.}

\paragraph{($\Leftarrow$)}  
Si $p = p_a$, alors toutes les votantes ont le même bulletin et approuvent la même personne.  
Donc, pour deux votes $a,a' \in \mathcal{A}$, on a $d_H(a,a') = 0$ car $a[i] = a'[i]$ et 
$\min d_H(a,a') = 0$. D'où $u_1^*(p) = 0$.

De même, on a $a_1 = a_2 = a$, donc :
\[
\sum \min \{d_H(a_1,a'), d_H(a_2,a')\} = u_2^*(p) = 0.
\]

Ainsi, pour $p = p_a$ :
\[
\varphi_{d_H}(p) = \frac{2}{n m}(0 - 0) = 0.
\]

\paragraph{($\Rightarrow$)} 

On suppose que $\varphi_{d_H}(p) = 0$.  
Par définition :
\[
\varphi_{d_H}(p) = \frac{2}{n m}\big(u_1^*(p) - u_2^*(p)\big),
\]
donc :
\[
\varphi_{d_H}(p) = 0 
\ \Leftrightarrow\ 
u_1^*(p) - u_2^*(p) = 0
\ \Leftrightarrow\ 
u_1^*(p) = u_2^*(p).
\]

Or, par définition :
\[
u_1^*(p) = \min_{a \in \mathcal{A}} \sum_{a' \in p} d_H(a, a')
\quad \text{et} \quad
u_2^*(p) = \min_{a_1, a_2 \in \mathcal{A}} 
\sum_{a' \in p} \{ d_H(a_1, a'), d_H(a_2, a') \}.
\]

On constate toujours que $u_2^*(p) \le u_1^*(p)$, car avec deux centres de cluster on peut toujours choisir $a_1=a_2$ et retrouver la même valeur que pour $u_1^*(p)$.  
Ainsi, l’égalité $u_2^*(p)=u_1^*(p)$ n’est possible que si l’ajout d’un second centroïde n’apporte strictement aucune amélioration du coût total, c’est‑à‑dire que tous les votes sont déjà identiques.

Autrement dit, cela n’est vrai que si :
\[
\forall a, a' \in p, \quad d_H(a, a') = 0,
\]
ce qui implique :
\[
a = a' \quad \forall a, a' \in p.
\]

Donc l’ensemble des bulletins contenus dans $p$ est réduit à un seul bulletin répété $n$ fois,  
autrement dit $p = p_a$ pour un certain $a \in \mathcal{A}$.


Ainsi, on a bien :
\[
\varphi_{d_H}(p) = 0 \ \Leftrightarrow\  p = p_a.
\]
\\	

\textbf{On montre que $\varphi_{d_H}(p) = 1$ ssi $p = p_{a, \bar{a}}$}

\paragraph{($\Leftarrow$)}  
Pour tout $\tilde{a} \in \mathcal{A}$ on a~: 
\[
\sum_{a \in p_{a,\bar{a}}} d_H(a,\tilde{a}) 
= \frac{n}{2} d_H(a,\tilde{a}) + \frac{n}{2} d_H(\bar{a}, \tilde{a}).
\]

%Si on choisit $\tilde{a} = a$, alors :
%\[
%\sum_{a \in p_{a,\bar{a}}} d_H(a,\tilde{a}) = \frac{n}{2} d_H(\bar{a}, a) 
%= \frac{n m}{2} 
%\quad \Rightarrow \quad 
%\text{distance } \leq \frac{n m}{2}.
%\]

Ainsi :

\begin{align*}
u_1^*(p_{a,\bar{a}}) 
&=
\min_{\tilde{a} \in \mathcal{A}} 
\sum_{a' \in p_{a,\bar{a}}} d_H(a', \tilde{a})\\
&=
\min_{\tilde{a} \in \mathcal{A}} 
\left(\frac{n}{2} d_H(a, \tilde{a}) + \frac{n}{2} d_H(\bar{a}, \tilde{a})\right)\\
&=
\frac{n}{2} \min_{\tilde{a} \in \mathcal{A}} \left( d_H(a, \tilde{a}) + d_H(\bar{a}, \tilde{a}) \right)
\end{align*}

Or, si $\tilde{a}$ diffère de $a$ sur $k$ positions, alors $\tilde{a}$ diffère de $\bar{a}$ sur $m - k$ positions. Donc~:
\[
d_H(a, \tilde{a}) + d_H(\bar{a}, \tilde{a}) = k + (m - k) = m.
\]

Ainsi :
\[
u_1^*(p_{a,\bar{a}}) = \frac{n}{2} \times m = \frac{n m}{2},
\]

Le calcul de $u_2^*(p_{a,\bar{a}})$ est plus simple~: en choisissant $a_1 = a$ et $a_2 = \bar{a}$, on trouve immédiatement les valeurs qui minimisent la somme et 
\[
u_2^*(p_{a,\bar{a}}) = 0.
\]

Ainsi~:

\[
\varphi(p_{a,\bar{a}}) 
= 
\frac{2}{nm} \left(u_1^*(p_{a,\bar{a}}) - u_2^*(p_{a,\bar{a}}) \right) 
= 
\frac{2}{nm} \left(\frac{nm}{2} - 0 \right) 
=
1.
\]

\paragraph{($\Rightarrow$)}

On a~:

\[
\varphi(p) = 1 
\Rightarrow \frac{2}{n m}(u_1^*(p) - u_2^*(p)) = 1 
\Rightarrow u_1^*(p) - u_2^*(p) = \frac{n m}{2}.
\]

On peut montrer que $u_1^*(p)$ est toujours bornée supérieurement par $\frac{nm}{2}$~: 
Si $a$ est un vote \emph{aléatoire}, où chaque bulletin vaut 1 avec probabilité $\frac12$, alors, en moyenne:

\begin{align*}
\text{E}_a \left( \sum_{a' \in p} d_H(a,a') \right)
&=
\sum_{a' \in p} \text{E}_a \left( d_H(a,a') \right) \quad \text{(linéarité de l'espérance)}\\
&=
\sum_{a' \in p} \frac{m}{2} \quad \text{(en moyenne, un vote aléatoire est à distance $m/2$ d'un vote)}\\
&=
\frac{nm}{2}
\end{align*}
Puisque la moyenne sur $a$ vaut $\frac{nm}{2}$, il existe au moins un bulletin pour lequel la somme est inférieure ou égale à $\frac{nm}{2}$. Par conséquent, on a toujours~:
\[
u_1^*(p) \le \frac{nm}{2}
\]

On déduit de $u_1^*(p) - u_2^*(p) = \frac{n m}{2}$ et de $u_1^*(p) \le \frac{nm}{2}$ que $u_2^*(p) \le 0$, et donc que 
\[
u_2^*(p) = 0
\]

puisque $u_2^*(p)$ est positif.

Or~:
\begin{align*}
u_2^*(p) = 0 
&\Leftrightarrow 
\sum \min \{d_H(a_1,a'), d_H(a_2,a')\} = 0\\
&\Leftrightarrow 
d_H(a_1,a') = 0 \text{ ou } d_H(a_2,a') = 0.
\end{align*}

Donc, pour tout $a' \in p$, on a $a' = a_1$ ou $a' = a_2$, il n’y a donc que deux types de bulletins. De plus 

\[
u_1^*(p) = \frac{n m}{2}.
\]

La méthode de raisonnement basée sur la moyenne effectuée plus haut, combinée au fait que $u_1^*(p) = \min_{\tilde{a}} \sum_{a \in p} d_H(\tilde{a},a) = \frac{n m}{2}$ permet de déduire que 
\[
\sum_{a \in p} d_H(\tilde{a},a) = \frac{n m}{2}
\quad
\text{pour tout $\tilde{a} \in A$}
\]
puisque si la valeur minimale est égale à la moyenne, toutes les valeurs sont égales à la moyenne.

Notons $k$ le nombre de bulletins $a_1$. On a $n-k$ bulletins $a_2$. On a

\[
\sum_{a \in p} d_H(a_1,a) = (n-k) d_H(a_1,a_2)
\quad\text{et}\quad
\sum_{a \in p} d_H(a_2,a) = k d_H(a_1,a_2)
\]

Les deux sommes valant, par ailleurs $\frac{nm}{2}$, on peut sommer les deux équations ci-dessus pour obtenir:

\[
d_H(a_1, a_2) = m
\]

Ce qui permet (enfin~!) de conclure que $a_2 = \bar{a_1}$, autrement dit, que $p = p_{a, \bar{a}}$.\\

Enfin, $\varphi_{d_H}(p) \in [0,1]$:
\begin{itemize}
	\item La positivité découle du fait que $u_1^*(p) \ge u_2^*(p)$, et
	\item comme on a vu qu'on a toujours $u_1^*(p) \le \frac{nm}{2}$, alors
\end{itemize}

\[
\varphi_{d_H}(p) 
\le
\frac{2}{nm}\left(\frac{nm}{2}-u_2^*(p)\right)
=
1-\frac{2}{nm}u_2^*(p)
\le
1
\]


Donc le premier axiome est vérifié.

\subsubsection{Axiome 2 : neutralité}

Pour que $\varphi_{d_H}(p) = \varphi_{d_H}(p^\pi)$ il suffit que $u_1^*(p) = u_1^*(p^\pi)$ et $u_2^*(p) = u_2^*(p^\pi)$. Montrons ces deux égalités.
\\

\begin{itemize}
  \item $u_1^*(p) = u_1^*(p^\pi)$
\end{itemize}

On a~:
\begin{align*}
u_1^*(p) 
&=
\min_a \sum_{a' \in p} d_H(a, a')\\
&=
\min_a \sum_{a' \in p} d_H(a^{\pi}, a'^{\pi})\\
&=
\min_a \sum_{a' \in p^{\pi}} d_H(a^{\pi}, a')\\
&=
\min_a \sum_{a' \in p^{\pi}} d_H(a, a')\\
&=
u_1^*(p^{\pi}) 
\end{align*}

\begin{itemize}
  \item $u_2^*(p) = u_2^*(p^\pi)$
\end{itemize}

On a~:
\begin{align*}
u_2^*(p)
&=
\min_{a_1,a_2}\sum_{a' \in p} \min \{ d_H(a_1, a'), d_H(a_2, a') \}\\
&=
\min_{a_1,a_2}\sum_{a' \in p} \min \{ d_H(a_1^{\pi}, a'^{\pi}), d_H(a_2^{\pi}, a'^{\pi}) \}\\
&=
\min_{a_1,a_2}\sum_{a' \in p^{\pi}} \min \{ d_H(a_1^{\pi}, a'), d_H(a_2^{\pi}, a') \}\\
&=
\min_{a_1,a_2}\sum_{a' \in p^{\pi}} \min \{ d_H(a_1, a'), d_H(a_2, a') \}\\
&=
u_2^*(p^{\pi})
\end{align*}

Comme on a montré que $u_1^*(p) = u_1^*(p^\pi)$ et $u_2^*(p) = u_2^*(p^\pi)$, on obtient :
\[
\frac{2}{n m}(u_1^*(p) - u_2^*(p))
= \frac{2}{n m}(u_1^*(p^\pi) - u_2^*(p^\pi))
\quad \Rightarrow \quad
\varphi_{d_H}(p) = \varphi_{d_H}(p^\pi).
\]
\\

Donc le second axiome est vérifié.


\subsubsection{Axiome 3 : anonymité}

\begin{itemize}
  \item $u_1^*(p) = u_1^*(\,^{\sigma}p)$
\end{itemize}

Pour tout $a \in \mathcal{A}$ :
\[
\sum_{a' \in p} d_H(a,a') 
= \sum_{a' \in {}^{\sigma}p} d_H(a,a'),
\]
car on ne fait que changer l’ordre des éléments de la somme (qui est finie).  
C’est donc aussi vrai pour le vote de consensus. Ainsi :
\[
u_1^*(p) = u_1^*(\,^{\sigma}p).
\]


\begin{itemize}
  \item $u_2^*(p) = u_2^*(\,^{\sigma}p)$
\end{itemize}

De manière analogue :
\[
\sum_{a' \in p} \min \{ d_H(a_1, a'), d_H(a_2, a') \}
= \sum_{a' \in {}^{\sigma}p} \min \{ d_H(a_1,a'), d_H(a_2,a') \}
\Rightarrow
u_2^*(p) = u_2^*(\,^{\sigma}p).
\]

Ainsi :
\[
\frac{2}{n m}(u_1^*(p) - u_2^*(p))
= \frac{2}{n m}(u_1^*(\,^{\sigma}p) - u_2^*(\,^{\sigma}p))
\Rightarrow
\varphi_{d_H}(p) = \varphi_{d_H}({}^{\sigma}p).
\]
\\

Donc le troisième axiome est vérifié.


\subsubsection{Axiome 4 : invariance par réplication}

\[
u_1^*(k p)
= \min_{\tilde{a} \in \mathcal{A}} \sum_{a' \in k p} d_H(\tilde{a}, a') 
= \min_{\tilde{a} \in \mathcal{A}} \sum_{a' \in p} k \, d_H(\tilde{a}, a')
= k u_1^*(p).
\]

De même :
\[
u_2^*(k p)
= \sum k \, \min_a \{ d_H(a_1,a'), d_H(a_2,a') \}
= k \sum \min_a \{ d_H(a_1,a'), d_H(a_2,a') \}
= k u_2^*(p).
\]

Ainsi :
\[
\varphi_{d_H}(k p)
= \frac{2}{(k n) m}(u_1^*(k p) - u_2^*(k p))
= \frac{2}{k n m}(k u_1^*(p) - k u_2^*(p))
= \frac{2}{n m}(u_1^*(p) - u_2^*(p))
= \varphi_{d_H}(p).
\]
\\

Note : On a bien $(kn)m$ au dénominateur de notre première égalité car on a $kn$ votes, grâce à la $k$-répétition de notre profil à $n$ votes.
\\

Donc le quatrième axiome est vérifié.


\subsection{Question 10}

Pour tout $a \in \mathcal{A}$~:

\[
\sum_{a' \in p} d_H(a,a') 
= \sum_{a' \in p} \sum_{i=1}^m (a[i] \oplus a'[i])
= \sum_{i=1}^m \sum_{a' \in p} (a[i] \oplus a'[i]).
\]

Le $a_i$ qui minimise $\sum_{a' \in p} (a[i] \oplus a'[i])$ est soit $0$, soit $1$, selon lequel des deux est le plus fréquent dans 
\[
(a'_{1}[i], a'_{2}[i], \ldots, a'_{n}[i]).
\]
On obtient donc $a$ en appliquant cette procédure pour chaque coordonnée.
\\

De manière rédigée :  
pour chaque coordonnée de $a$, on regarde quel chiffre est majoritaire à cette position dans les votes.  
S’il s’agit de $0$, alors notre consensus aura un $0$ à cette coordonnée (car cela représentera la décision majoritaire du profil).  
Sinon, ce sera un $1$.  
En cas d’égalité, le choix importe peu.


\subsection{Question 11}

On sait que :
\[
d_S(\succ, \succ_{v_k}) 
= \sum_{c_j \in C} |r_\succ(c_j) - r_{\succ_{v_k}}(c_j)|.
\]

On cherche à minimiser la somme de ces distances sur tous les profils~:
\[
\sum_{v \in p} \sum_{c_j \in C} 
|r_\succ(c_j) - r_{\succ_{v_k}}(c_j)|
= 
\sum_{c_j \in C} \sum_{v \in p} 
|r_\succ(c_j) - r_{\succ_{v_k}}(c_j)|.
\]

Pour chaque candidate, on cherche un rang unique tel que la somme totale des écarts soit minimale.

En utilisant le principe du graphe biparti (c’est-à-dire qu’on peut découper le graphe en deux ensembles de sommets tels que deux sommets d’un même ensemble ne sont jamais adjacents),  
on pose :

\begin{itemize}
	\item le premier ensemble de sommets :  
  \( U \) = les \( m \) candidates ;
  \item le second ensemble de sommets :  
  \( V \) = les \( m \) positions possibles.
\end{itemize}

On commence par dire que chaque candidate est connectée à toutes les positions possibles.  
On associe un poids à chaque arête correspondant au coût total si la candidate en question est placée au rang \( k \) :

\[
w_{j,k} = \sum_{v \in p} |\, k - r_{\succ_v}(c_j)|.
\]

Il suffit de trouver la combinaison qui minimise la somme finale.  
Cette combinaison doit forcer notre graphe à être en couplage parfait  
(c’est-à-dire qu’aucune candidate n’est associée à deux positions différentes et qu’aucune position n’est associée à deux candidates différentes).  
On retrouve ainsi la bijection recherchée.

Cette recherche de combinaison se fait à l’aide du module \texttt{scipy} sous Python  
(voir la question 12 pour l’implémentation).

\subsection{Question 15}
Pour étudier empiriquement le comportement des mesure de polarisation $\varphi_{d_H}$ et $\varphi_{d_S}$, il faut simuler en faisant varier le paramètre de polarisation $p$ du modèle de génération des profils.

Les résultats sont représentés graphiquement en fonction de $p$ et on observe que la polarisation est minimale lorsque $p = 0$ et maximale autour de $p = 1$, ce qui correspond à une situation où la population est divisée en deux groupes de taille comparable adoptant des préférences opposées. 

Ces résultats sont ceux attendus car ils montrent que $\varphi_{d_H}$ (respectivement $\varphi_{d_S}$) est bien une mesure de polarisation pour le cas des votes par approbation (respectivement, votes par ordres totaux).


\section{Utilisation des IAs}

\subsection{BENAMEUR Féryel}
\textbf{Prompt 1 :} je suis en train de montrer l'axiome 1 et comme il y a un ssi je dois faire le sens ou si phi(p) = 0 alors p = pa et de même pour l'autre. dis moi si je pars bien (cf pièce jointe).

\textbf{Réponse :} Il me dit que ce que je fais est bien mais trop compliqué. Il m’a donné une méthode pour reprendre ma démonstration et mes calculs de manière simplifiée.

\textbf{Critique :} Son avis ne m’a pas été d’une grande utilité. Finalement, je n’ai rien changé de ce que j’avais fait au préalable.
\\

\textbf{Prompt 2 :} j’ai testé ça pour le sens où phi(p) = 1 alors p = pa,ā. Je n’arrive pas à continuer ma preuve. Reprends l’idée de ma démonstration pour phi(p) = 0 et applique la à notre situation.


\textbf{Réponse :} explique, étape par étape, la procédure à suivre pour résoudre le problème. 

\textbf{Critique :} J’ai globalement repris l’idée de ses quatre étapes et les ai un peu reformulé à ma manière. J’ai bien aimé le fait qu’il découpe la démonstration en 4 étapes qui se suivent afin de me guider. De plus, sa preuve s’inspirait beaucoup de la mienne pour phi(p) = 0 donc c’était plus facile et clair pour moi de le comprendre et de rédiger à ma façon.
\\

\textbf{Prompt 3 :} Comment transformer un notebook jupyter en fichier python éxecutable depuis le terminal ?

\textbf{Réponse :} Il me dit d’utiliser la commande “python3 -m nbconvert --to script projet.ipynb”

\textbf{Critique :} ça fonctionne parfaitement

\subsection{BOUCHERON Lucie}

\textbf{Prompt 1 :} Voici la définition de $\varphi^2$ : $\varphi^2(p) = \sum_{\{c_k,c_l\} \in C^2} \frac{n - d_{c_k,c_l}(p)}{n \binom{m}{2}}$. Je cherche à simplifier son implémentation. 
Est-ce qu’il existe une écriture plus compacte, plus computationnelle qui permettrait une implémentation plus simple en Python ?

\textbf{Réponse :}  L’IA propose de réécrire la formule de $\varphi^2$ d'une certaine façon, en expliquant que cela permet de sommer directement les valeurs d’une matrice contenant les différences absolues, ce qui permet d’utiliser une seule somme au lieu d’une double boucle sur les paires dans l’implémentation.

\textbf{Critique :} Cette réponse nous a été utile pour simplifier l’écriture du code mais l’IA ne justifie pas en détail le passage d’une forme à l’autre, donc il a fallu vérifier par nous même l’équivalence.
\\

\textbf{Prompt 2 :} Je veux échanger deux valeurs dans un tableau sur python, sans que les autres valeurs de mon tableau ne bougent, quelle commande je peux utiliser ?

\textbf{Réponse :}  Elle propose np.ramdom.choice

\textbf{Critique :} Cette réponse direct et efficace a été immédiatement utilisée dans mon code et a fonctionné.
\\

\textbf{Prompt 3 :} Dans mon modèle de génération de votes par approbation et par ordre, quand je teste mes algorithmes, j’obtiens des profils assez homogènes, même en faisant varier le paramètre de polarisation. Sans modifier la structure principale du modèle, ni le fond de mon algorithme donne moi une façon d’avoir plus de variabilité dans les groupes pour rendre les profils plus réalistes. 

\textbf{Réponse :} ajout du bruit 

\textbf{Critique :} Cette réponse a été très utile pour la suite du projet, et l’implémentation est très simple et efficace. 

\subsection{BAIGNERES Clara}

\textbf{Prompt 1} : Je suis en train de rédiger le rapport pour un projet en fondement mathématique pour l'aide à la décision. Je te donne le sujet en pièce jointe. Je fais ce rapport en \LaTeX. Je vais te donner des morceaux de codes \LaTeX parties par parties. Je veux que tu les mettes en page, c'est à dire que la présentation soit belle et lisible, surtout pour les calculs. Je veux un style purement académique et ne modifie rien à ma formulation.

\textbf{Réponse} : pour chaque blocs que je lui ai envoyé, il m'a rendu, mot à mot, le même texte \LaTeX, mais avec une mise en page aérée, bien plus agréable. 

\textbf{Critique :} cela m'a été grandement utile, j'ai économisé beaucoup de temps simplement pour la mise en page.
\\

\textbf{Prompt 2 :} prouve moi l'axiome 1 de l'exercice 9 dans le cas où phi(p) = 1 

\textbf{Réponse :} il me propose une preuve telle qu'un de mes centroïdes est forcément un vote de mon profil, et non pas un consensus. 

\textbf{Critique :} je ne comprenais pas la preuve de départ de Féryel donc j'ai demandé à Mistral une autre version. Sa réponse ne m'a été que partiellement utile car comme dit juste au dessus, l'IA ne comprenait pas le principe de consensus, et m'attribuait $u_1^*$ à un vote de mon profil. Cependant, il m'a donné l'idée de la preuve pour $u_1^* \leq \frac{nm}{2}$, ce qui m'a permis de compléter la réponse de Féryel.  

\section{Sources}


\begin{itemize}
	\item \underline{\href{http://pingouindesalpes.com/hugo/post/markdown/}{Notation markdown pour le ReadMe}.}
	\item \underline{\href{https://share.miple.co/content/6IggazU5m0tIG
}{Démonstration de l’inégalité triangulaire pour la valeur absolue}.}
\end{itemize}

































\end{document}
